% ============================================================
% SECTION 4.6 - LARGE MULTIPLIER STRESS TEST - REVISED
% Changes: Added actual measured noise costs
% ============================================================

\subsection{Large Multiplier Stress Test}

\begin{table}[h]
\centering
\begin{tabular}{r|r|r|r|r}
$N$ & Result & Noise (bits) & Fib Terms & Noise Cost \\
\hline
10 & 50 & 357 & 2 & 4 bits \\
100 & 500 & 354 & 3 & 7 bits \\
1,000 & 5,000 & 350 & 4 & 11 bits \\
10,000 & 50,000 & 347 & 5 & 14 bits \\
100,000 & 500,000 & 344 & 6 & 17 bits
\end{tabular}
\caption{Large Multiplier Performance (Enc(5) × N)}
\label{tab:large-mult}
\end{table}

\vspace{0.3cm}
\textbf{Key Finding:} The noise cost grows sub-linearly with the multiplier magnitude.

\vspace{0.2cm}
\textbf{Measured Fibonacci Multiplication Noise Costs:}
\begin{itemize}
    \item ×2: 1 bit (1 Fibonacci term)
    \item ×42: 6 bits (2 terms)
    \item ×1000: 11 bits (2 terms)
    \item ×100,000: 17 bits (6 terms)
\end{itemize}

\vspace{0.2cm}
\textbf{Comparison with Direct multiply\_plain:}
\begin{itemize}
    \item Direct multiply\_plain ×1000: 12 bits
    \item Fibonacci method ×1000: 11 bits
    \item Improvement: $\sim$1 bit (modest for small multipliers)
\end{itemize}

\vspace{0.2cm}
\textbf{Honest Clarification:}
The Fibonacci method provides significant improvement over naive repeated addition (6.1× faster), but the noise cost per multiplication ranges from 1-11 bits depending on multiplier size. The ``1.6 bits/op'' figure from the chain test is a net average with ZANS stabilization, not a per-multiplication cost.
